\section{AHSME 1971}
        \begin{enumerate}
        	\item (\textit{AHSME 1971}) The number of digits in the number $N=2^{12}\times 5^8$ is


 $\textbf{(A) }9\qquad \textbf{(B) }10\qquad \textbf{(C) }11\qquad \textbf{(D) }12\qquad  \textbf{(E) }20$


 
 \textbf{Answer: }B



\textbf{Solution}

$N=2^{12}\cdot5^8=2^4\cdot2^8\cdot5^8=2^4\cdot(2\cdot5)^8=2^4\cdot10^8=16\cdot100000000$, which has $2+8=10$ digits, so the answer is $\boxed{\textbf{(B) }10}$.

 


	\item (\textit{AHSME 1971}) If $b$ men take $c$ days to lay $f$ bricks, then the number of days it will take $c$ men working at the same rate to lay $b$ bricks, is


 $\textbf{(A) }fb^2\qquad \textbf{(B) }b/f^2\qquad \textbf{(C) }f^2/b\qquad \textbf{(D) }b^2/f\qquad  \textbf{(E) }f/b^2$


 
 \textbf{Answer: }D



\textbf{Solution}

We can use a modified version of the equation $\text{Distance} = \text{Rate} \times {\text{Time}}$, which is $\text{Work Done} = \text{Rate of Work} \times{ \text{Time Worked}}$. In this case, the work done is the number of bricks laid, the rate of work is the number of men working, and the time worked is the number of days. With these definitions, we see from the first equation that $f = b \cdot c$. If we let $X$ be the number of days it will take $c$ men working at the same rate to $b$ bricks, then we have the equation $b = c \cdot x$. So, $x = \frac{b}{c}$. The first equation says that $\frac{f}{c^2} = \frac{b}{c}$, which leads to $x = \frac{f}{c^2}$. This doesn't match any of our answer choices though, so we have to fiddle around a bit before we realize that $c = \frac{f}{b}$, a substitution we can make to see that $x = \frac{b^2}{f}$. Thus, the answer is $\boxed{\textbf{(D) }b^2/f}$.

 


	\item (\textit{AHSME 1971}) If the point $(x,-4)$ lies on the straight line joining the points $(0,8)$ and $(-4,0)$ in the $xy$-plane, then $x$ is equal to


 $\textbf{(A) }-2\qquad \textbf{(B) }2\qquad \textbf{(C) }-8\qquad \textbf{(D) }6\qquad  \textbf{(E) }-6$


 
 \textbf{Answer: }E



\textbf{Solution}

Since $(x,-4)$ is on the line $(0,8)$ and $(-4,0)$, the slope between the latter two is the same as the slope between the former two, so $\frac{-4-8}{x-0}=\frac{8-0}{0-(-4)}$. Solving for $x$, we get $x=-6$, so the answer is $\boxed{\textbf{(E) }-6}$.

 


	\item (\textit{AHSME 1971}) After simple interest for two months at $5$\% per annum was credited, a Boy Scout Troop had a total 
of $\textdollar 255.31$ in the Council Treasury. The interest credited was a number of dollars plus the following number of cents


 $\textbf{(A) }11\qquad \textbf{(B) }12\qquad \textbf{(C) }13\qquad \textbf{(D) }21\qquad  \textbf{(E) }31$


 
 \textbf{Answer: }A



\textbf{Solution}

We can use the formula $A = P(1+rt)$ for simple interest, where $P$ is the principal (initial amount), $A$ is the total, $r$ is the interest rate per year, and $t$ is the time in years.

 We are solving for $P,$ so we can rearrange the equation:
 \[P = \frac{A}{1+rt}\]Plug in all the necessary numbers:
 \[P = \frac{255.31}{1 + (0.05)(2/12)} = \textdollar 253.21\]So, the interest is 
 \[P*(0.05)(2/12) = \textdollar 2.11\]

 The answer is $\boxed{\textbf{(A) }11}$

 -edited by Happypuma46

 If there are any mistakes, feel free to edit so that it is correct.

 


	\item (\textit{AHSME 1971}) Points $A,B,Q,D$, and $C$ lie on the circle shown and the measures of arcs $\widehat{BQ}$ and $\widehat{QD}$ 
are $42^\circ$ and $38^\circ$ respectively. The sum of the measures of angles $P$ and $Q$ is


 $\textbf{(A) }80^\circ\qquad \textbf{(B) }62^\circ\qquad \textbf{(C) }40^\circ\qquad \textbf{(D) }46^\circ\qquad  \textbf{(E) }\text{None of these}$


 
                    \begin{figure}[H]
                        centering
                        \includegraphics[width=0.5\textwidth]{Problem/AHSME_I_1971_5_Problem_Figure_1.png}
                    \end{figure}
                    


 
 \textbf{Answer: }C



\textbf{Solution 1}

We see that the measure of $P$ equals $(\widehat{BD}-\widehat{AC})/2$, and that the measure of $Q$ equals $\widehat{AC}/2$.Since $\widehat{BD} = \widehat{BQ} + \widehat{QD} = 42^{\circ} + 38^{\circ} = 80^{\circ}$, the sum of the measures of $P$ and $Q$ is $\widehat{BD}/2 = 80^{\circ}/2 = 40^{\circ} \Longrightarrow \boxed{\textbf{(C) }40^{\circ}}$.

 



\textbf{Solution 2}

Arcs are measured by the angle measures of their corresponding \href{https://artofproblemsolving.com/wiki/index.php?title=Central\_angle}{central angles}. Thus, the \href{https://artofproblemsolving.com/wiki/index.php?title=Inscribed\_angle}{inscribed angle} $\measuredangle BAQ = \frac{42^{\circ}}2 = 21^{\circ}$, and, likewise, $\measuredangle QCD = \frac{38^{\circ}}2 = 19^{\circ}$. Thus, by \href{https://artofproblemsolving.com/wiki/index.php?title=Supplementary}{supplementary} angles, $\measuredangle PAQ = 180^{\circ} - \measuredangle BAQ = 159^{\circ}$, and $\measuredangle PCQ = 180^{\circ} - \measuredangle QCD = 161^{\circ}$. Because the sum of the \href{https://artofproblemsolving.com/wiki/index.php?title=Interior\_angle}{interior angle} measures of a quadrilateral add to $360^{\circ}$, we see that $\measuredangle P + \measuredangle Q = 360^{\circ} - \measuredangle PAQ - \measuredangle PCQ = 360^{\circ} - 159^{\circ} - 161^{\circ} = 40^{\circ}$. Thus, our answer is $\boxed{\textbf{(C) }40^{\circ}}$.

 


	\item (\textit{AHSME 1971}) Let $\ast$ be the symbol denoting the binary operation on the set $S$ of all non-zero real numbers as follows: 
For any two numbers $a$ and $b$ of $S$, $a\ast b=2ab$. Then the one of the following statements which is not true, is


 $\textbf{(A) }\ast\text{ is commutative over }S \qquad \textbf{(B) }\ast\text{ is associative over }S\qquad \\ \textbf{(C) }\frac{1}{2}\text{ is an identity element for }\ast\text{ in }S\qquad  \textbf{(D) }\text{Every element of }S\text{ has an inverse for }\ast\qquad  \textbf{(E) }\dfrac{1}{2a}\text{ is an inverse for }\ast\text{ of the element }a\text{ of }S$


 
 \textbf{Answer: }E



\textbf{Solution}

$\textbf{(A) }\ast\text{ is commutative over }S$
 \[a \ast b = b \ast a = 2ab\]Statement A is true.

 


 $\textbf{(B) }\ast\text{ is associative over }S$
 \[a \ast (b \ast c) = a \ast 2bc = 2a2bc = 4abc\]
 \[(a \ast b) \ast c = 2ab \ast c = 2(2ab)(c) = 4abc\]Statement B is true.

 
$\textbf{(C) }\frac{1}{2}\text{ is an identity element for }\ast\text{ in }S\qquad$
 \[a \ast \frac{1}{2} = a\]Statement C is true.

 $\textbf{(D) }\text{From the previous answer choice, we know }\frac{1}{2}\text{ is the identity element for }\ast\text{ in }S.\qquad$$\text{For inverses to exist, they must evaluate to the identity under the operator }\ast.\text{ Thus, } a \ast b = \frac{1}{2}\text{, which leads to } b = \frac{1}{4a}$

 Statement D is true.

 
$\textbf{(E) }\dfrac{1}{2a}\text{ is an inverse for }\ast\text{ of the element }a\text{ of }S$
 \[\frac{1}{2a} \ast a = 1\]

 Since the identity is $\frac{1}{2}$, not 1, we can see that statement E is false.

 Hence, the answer is $\boxed{\textbf{(E)}}$.

 -edited by coolmath34-fixed by DoctorSeventeen

 


	\item (\textit{AHSME 1971}) $2^{-(2k+1)}-2^{-(2k-1)}+2^{-2k}$ is equal to 


 $\textbf{(A) }2^{-2k}\qquad \textbf{(B) }2^{-(2k-1)}\qquad \textbf{(C) }-2^{-(2k+1)}\qquad \textbf{(D) }0\qquad  \textbf{(E) }2$


 
 \textbf{Answer: }C



\textbf{Solution}

By using the properties of \href{https://artofproblemsolving.com/wiki/index.php?title=Exponentiation\#Basic\_Properties}{exponents}, we can simplify the given expression as follows to obtain our answer:\begin{align*}2^{-(2k+1)} - 2^{-(2k-1)} + 2^{-2k} &= 2^{-2k-1} - 2^{-2k+1} + 2^{-2k} \\&= \frac{2^{-2k}}2 - 2\cdot2^{-2k} + 2^{-2k} \\&= 2^{-2k}(\frac12 - 2 + 1) \\&= 2^{-2k}(-\frac12) \\&= -\frac{2^{-2k}}2 \\&= -2^{-2k-1} \\&= \boxed{\textbf{(C) }-2^{-(2k+1)}}.\end{align*}


 


	\item (\textit{AHSME 1971}) The solution set of $6x^2+5x<4$ is the set of all values of $x$ such that


 $\textbf{(A) }\textstyle -2<x<1\qquad \textbf{(B) }-\frac{4}{3}<x<\frac{1}{2}\qquad \textbf{(C) }-\frac{1}{2}<x<\frac{4}{3}\qquad \\ \textbf{(D) }x<\textstyle\frac{1}{2}\text{ or }x>-\frac{4}{3}\qquad \textbf{(E) }x<-\frac{4}{3}\text{ or }x>\frac{1}{2}$


 
 \textbf{Answer: }B



\textbf{Solution}

We are solving the inequality $6x^2 + 5x - 4  < 0.$  This can be factored as
 \[(2x-1)(3x+4) < 0\]

 The graph of this inequality is a parabola facing upwards, so the \href{https://artofproblemsolving.com/wiki/index.php?title=Interval}{interval} between the roots satisfies the equation. This interval, $(-\frac{4}{3}, \frac{1}{2})$, is answer $\boxed{\textbf{(B)}}$.

 -edited by coolmath34

 


	\item (\textit{AHSME 1971}) An uncrossed belt is fitted without slack around two circular pulleys with radii of $14$ inches and $4$ inches. 
If the distance between the points of contact of the belt with the pulleys is $24$ inches, then the distance 
between the centers of the pulleys in inches is


 $\textbf{(A) }24\qquad \textbf{(B) }2\sqrt{119}\qquad \textbf{(C) }25\qquad \textbf{(D) }26\qquad  \textbf{(E) }4\sqrt{35}$


 
 \textbf{Answer: }D



\textbf{Solution}


                    \begin{figure}[H]
                        centering
                        \includegraphics[width=0.5\textwidth]{Solution/AHSME_I_1971_9_Solution_Figure_1.png}
                    \end{figure}
                    

 Let the center of the smaller circle be $A$ and the center of the larger circle be $B$. Draw $\overline{AC}$ and $\overline{BD}$, as in the diagram. We know that $\overline{AC}$ and $\overline{BD}$ are perpendicular to $\overline{CD}$ because $\overline{CD}$ is \href{https://artofproblemsolving.com/wiki/index.php?title=Tangent\_(geometry)\#Tangents\_to\_Circles}{tangent} to both circles. From the problem, we know $CD=24$.

 Draw $\overline{AE} \perp \overline{BD}$ with $E$ on $\overline{BD}.$ $ABDC$ is a \href{https://artofproblemsolving.com/wiki/index.php?title=Rectangle}{rectangle}, so, knowing the radii of the circles, we see that $BE = BD - ED = BD - AC = 14 - 4 = 10.$

 Points $A, E,$ and $B$ form a right triangle with legs $10$ and $24.$ We are looking for $AB,$ which is $\sqrt{10^2 + 24^2} = 26$ by the \href{https://artofproblemsolving.com/wiki/index.php?title=Pythagorean\_Theorem}{Pythagorean Theorem}.

 Thus, our answer is $\boxed{\textbf{(D) }26}$.

 -edited by coolmath34

 


	\item (\textit{AHSME 1971}) Each of a group of $50$ girls is blonde or brunette and is blue eyed or brown eyed. If $14$ are blue-eyed blondes, 
$31$ are brunettes, and $18$ are brown-eyed, then the number of brown-eyed brunettes is


 $\textbf{(A) }5\qquad \textbf{(B) }7\qquad \textbf{(C) }9\qquad \textbf{(D) }11\qquad  \textbf{(E) }13$


 
 \textbf{Answer: }E



\textbf{Solution}

There are $31$ brunettes, so there are $50-31=19$ blondes. We know there are $14$ blue-eyed blondes, so there are $19-14=5$ brown-eyed blondes.

 Next, we know that $18$ people are brown-eyed, so there are $18-5=13$ brown-eyed brunettes.

 The answer is $\boxed{\textbf{(E) }13}.$

 -edited by coolmath34

 


	\item (\textit{AHSME 1971}) The numeral $47$ in base a represents the same number as $74$ in base $b$. Assuming that both bases are positive 
integers, the least possible value of $a+b$ written as a Roman numeral, is


 $\textbf{(A) }\mathrm{XIII}\qquad \textbf{(B) }\mathrm{XV}\qquad \textbf{(C) }\mathrm{XXI}\qquad \textbf{(D) }\mathrm{XXIV}\qquad \textbf{(E) }\mathrm{XVI}$


 
 \textbf{Answer: }D



\textbf{Solution}

$4a+7=7b+4\implies 7b=4a+3\implies b=\frac{4a+3}{7}$. $4a+3\equiv 0\implies 4a\equiv 4\implies a\equiv1 \pmod{7}$. The smallest possible value of $a$ is $8$. Then, $b=\frac{4\cdot8+3}{7}=\frac{35}{7}=5$. However, the digit $7$ is not valid in base $5$, so we have to try a larger value. $a=8+7=15$, gives a value of $\frac{15*4+3}{7}=9$, for $b$, which is valid. 

 $15+9=24$, which is $\mathrm{XXIV}$ as a roman numeral, and thus the answer is $\boxed{\textbf{(D) }\mathrm{XXIV}}.$

 


	\item (\textit{AHSME 1971}) For each integer $N>1$, there is a mathematical system in which two or more positive integers are defined 
to be congruent if they leave the same non-negative remainder when divided by $N.$ If $69,90$, and $125$ are 
congruent in one such system, then in that same system, $81$ is congruent to


 $\textbf{(A) }3\qquad \textbf{(B) }4\qquad \textbf{(C) }5\qquad \textbf{(D) }7\qquad  \textbf{(E) }8$


 
 \textbf{Answer: }B



\textbf{Solution}

The "mathematical system" being alluded to is \href{https://artofproblemsolving.com/wiki/index.php?title=Modulo}{modulo} $N$, so we shall be using such notation for convenience. From the problem, we know that $69 \equiv 90 \equiv 125 \mod N$, so the difference between each of these numbers must be a multiple of $N$. $90-69=21=3\cdot7$, and $125-90=35=5\cdot7$. The only common factors between these two numbers are $1$ and $7$, but we know from the problem that $N\neq1$. Thus, $N=7$. Now, $81 \equiv 77+4 \equiv 4 \mod 7$, so our answer is $\boxed{\textbf{(B) }4}$.

 


	\item (\textit{AHSME 1971}) If $(1.0025)^{10}$ is evaluated correct to $5$ decimal places, then the digit in the fifth decimal place is


 $\textbf{(A) }0\qquad \textbf{(B) }1\qquad \textbf{(C) }2\qquad \textbf{(D) }5\qquad  \textbf{(E) }8$


 
 \textbf{Answer: }E



\textbf{Solution}

We can rewrite $(1.0025)^{10}$ as $(1+\frac{25}{10^4})^{10}$. By the \href{https://artofproblemsolving.com/wiki/index.php?title=Binomial\_Theorem}{Binomial Theorem}, we know that this expression equals $1+ \tbinom{10}1\cdot\frac{25}{10^4}+\tbinom{10}2\cdot\frac{25^2}{10^8}+\tbinom{10}3\cdot\frac{25^3}{10^{12}}+\cdots+(\frac{25}{10^4})^{10}$. We are looking for the $\frac1{10^5}$ term (the fifth decimal digit), so we can disregard the first two terms, because they will not affect our final result. Thus, we are left with $\tbinom{10}2\cdot\frac{25^2}{10^8}+\tbinom{10}3\cdot\frac{25^3}{10^{12}}+\cdots = \frac{45\cdot25^2}{10^8}+\frac{120\cdot25^3}{10^{12}}+\cdots = \frac{28,125}{10^8}+\frac{1,875,000}{10^{12}}+\cdots$. The first term here has a $\frac1{10^5}$ term of $8$, and the second term does not contribute to the fifth decimal place, because it does not have enough digits in the numerator. So, we can disregard the subsequent terms, because they will not have enough digits either. Thus, the fifth digit of the expansion is $\boxed{\textbf{(E) }8}$.

 


	\item (\textit{AHSME 1971}) The number $(2^{48}-1)$ is exactly divisible by two numbers between $60$ and $70$. These numbers are


 $\textbf{(A) }61,63\qquad \textbf{(B) }61,65\qquad \textbf{(C) }63,65\qquad \textbf{(D) }63,67\qquad  \textbf{(E) }67,69$


 
 \textbf{Answer: }C



\textbf{Solution}

Factor repeatedly with \href{https://artofproblemsolving.com/wiki/index.php?title=Difference\_of\_squares}{difference of squares}.
 \[2^{48}-1 = (2^{24}+1)(2^{12}+1)(2^{6}+1)(2^6-1)\]

 We only care about two terms: $2^6+1$ and $2^6-1$. These simplify to $65$ and $63$.

 Thus, our answer is $\boxed{\textbf{(C) }63,65}.$

 -edited by coolmath34

 


	\item (\textit{AHSME 1971}) An aquarium on a level table has rectangular faces and is $10$ inches wide and $8$ inches high. When it was tilted, 
the water in it covered an $8\times 10$ end but only $\frac{3}{4}$ of the rectangular bottom. The depth of the 
water when the bottom was again made level, was


 $\textbf{(A) }2\textstyle{\frac{1}{2}}"\qquad \textbf{(B) }3"\qquad \textbf{(C) }3\textstyle{\frac{1}{4}}"\qquad \textbf{(D) }3\textstyle{\frac{1}{2}}"\qquad \textbf{(E) }4"$


 
 \textbf{Answer: }B



\textbf{Solution}

Let the third dimension of the aquarium be $x$, and let the height of the water when the aquarium is level be $h$. When the aquarium is tilted, the water forms a triangular prism. The triangular faces have height $8$ and, from the problem, base $\frac34x$. Thus, they have area $\frac12\cdot8\cdot\frac34x=3x$, so, because the prism has length $10$, the volume of the water is $10\cdot3x=30x$. When the tank is level, the water forms a rectangular prism with volume $h\cdot10\cdot x=10hx$. Because the amount of water in the tank is conserved, we can equate these two expressions for the volume of the water. Thus, $10hx=30x$, so $h=\boxed{\textbf{(B) }3"}$.

 


	\item (\textit{AHSME 1971}) After finding the average of $35$ scores, a student carelessly included the average with the $35$ scores and found the 
average of these $36$ numbers. The ratio of the second average to the true average was


 $\textbf{(A) }1:1\qquad \textbf{(B) }35:36\qquad \textbf{(C) }36:35\qquad \textbf{(D) }2:1\qquad  \textbf{(E) }\text{None of these}$


 
 \textbf{Answer: }A



\textbf{Solution}

Assume the $35$ scores are the first $35$ natural numbers: $1, 2, 3, \dots 35.$

 The average of the scores is $\frac{\frac{(35)(36)}{2}}{35} = 18.$ 

 If we add $18$ to the first $35$ numbers, the new average is $\frac{648}{36} = 18$ still.

 The two averages are the same, therefore the answer is $\boxed{\textbf{(A) }1:1}.$

 -edited by coolmath34

 


	\item (\textit{AHSME 1971}) A circular disk is divided by $2n$ equally spaced radii($n>0$) and one secant line. The maximum number of non-overlapping 
areas into which the disk can be divided is


 $\textbf{(A) }2n+1\qquad \textbf{(B) }2n+2\qquad \textbf{(C) }3n-1\qquad \textbf{(D) }3n\qquad  \textbf{(E) }3n+1$


 
 \textbf{Answer: }E



\textbf{Solution 1}

Note that, before adding the secant line, we have $2n$ regions, and, starting from the edge of the disk, every time the secant line intersects a radius or the other end of the disk, a new region is created. Thus, we desire to maximize the number of radii that the secant line passes through.

 The angle between any two adjacent radii must be $\frac{360^{\circ}}{2n}=\frac{180^{\circ}}n$. A given secant cannot pass through more than $180^{\circ}$ of the radii, and passing through exactly $180^{\circ}$ would imply that the secant must go through the center of the circle, and so it would not produce any new regions. Thus, to maximize the number of regions, the secant line must pass through $n$ radii (spanning $\frac{180}n(n-1)$ degrees) and thereby create $n+1$ new regions (one for each radius and one for exiting the disk). Thus, we have a maximum of $2n+n+1=\boxed{\textbf{(E) }3n+1}$ regions.

 



\textbf{Solution 2}

We can draw the cases for small values of $n.$
 \[n = 0 \rightarrow \text{areas} = 1\]
 \[n = 1 \rightarrow \text{areas} = 4\]
 \[n = 2 \rightarrow \text{areas} = 7\]
 \[n = 3 \rightarrow \text{areas} = 10\]It seems that for $2n$ radii, there are $3n+1$ distinct areas. The secant line must pass through $n$ radii for this to occur.

 The answer is $\boxed{\textbf{(E) } 3n+1}.$

 -edited by coolmath34

 


	\item (\textit{AHSME 1971}) The current in a river is flowing steadily at $3$ miles per hour. A motor boat which travels at a constant rate in still 
water goes downstream $4$ miles and then returns to its starting point. The trip takes one hour, excluding the time spent 
in turning the boat around. The ratio of the downstream to the upstream rate is 


 $\textbf{(A) }4:3\qquad \textbf{(B) }3:2\qquad \textbf{(C) }5:3\qquad \textbf{(D) }2:1\qquad  \textbf{(E) }5:2$


 
 \textbf{Answer: }D



\textbf{Solution}

Let the boat travel over still water at rate $r$. Then, it travels downstream at rate $r+3$ and upstream at rate $r-3$. The distance travelled downstream can be expressed as $(r+3)t$, where $t$ is the time taken to travel downstream. Because the total time taken is $1$ hour, the distance travelled upstream can be represented as $(r-3)(1-t)$. Because both distances are the same ($4$ miles), we can equate the two expressions to solve for $t$ in terms of $r$:\begin{align*}(r+3)t &= (r-3)(1-t) \\rt+3t &= r-rt-3+3t \\2rt &= r-3 \\t &= \frac{r-3}{2r} \\\end{align*}
Recalling that the distance travelled downstream, $(r+3)t$, is $4$ miles, we can substitute the above expression for $t$ to solve for $r$:\begin{align*}(r+3)t &= 4 \\(r+3)\frac{r-3}{2r} &= 4 \\r^2-9 &= 8r \\r^2-8r-9 &= 0 \\(r-9)(r+1) &= 0 \\\end{align*}
Because $r>0$, $r=9$. Thus, the ratio of the downstream rate $r+3$ to the upstream rate $r-3$ is $\frac{9+3}{9-3}=\frac{12}6=\boxed{\textbf{(D) }2:1}$.

 


	\item (\textit{AHSME 1971}) If the line $y=mx+1$ intersects the ellipse $x^2+4y^2=1$ exactly once, then the value of $m^2$ is


 $\textbf{(A) }\textstyle\frac{1}{2}\qquad \textbf{(B) }\frac{2}{3}\qquad \textbf{(C) }\frac{3}{4}\qquad \textbf{(D) }\frac{4}{5}\qquad \textbf{(E) }\frac{5}{6}$


 
 \textbf{Answer: }C



\textbf{Solution}

Plug in $y=mx+1$ into the ellipse's equation to find the intersection points:
 \[x^2 + 4(mx+1)^2 = 1\]After simplifying, we have a quadratic in $x$:
 \[(4m^2 + 1)x^2 + 8mx  +3 = 0\]

 Because there is only one intersection point, then the quadratic has only one solution. This can only happen when the \href{https://artofproblemsolving.com/wiki/index.php?title=Discriminant}{discriminant} is 0.
 \[\Delta = b^2 - 4ac = (8m)^2 - (4)(1+4m^2)(3) = 0\]

 Solving, we find $m^2 = \boxed{\textbf{(C) }\frac34}.$

 -edited by coolmath34

 


	\item (\textit{AHSME 1971}) The sum of the squares of the roots of the equation $x^2+2hx=3$ is $10$. The absolute value of $h$ is equal to


 $\textbf{(A) }-1\qquad \textbf{(B) }\textstyle\frac{1}{2}\qquad \textbf{(C) }\textstyle\frac{3}{2}\qquad \textbf{(D) }2\qquad \textbf{(E) }\text{None of these}$


 
 \textbf{Answer: }E



\textbf{Solution 1}

We can rewrite the equation as $x^2 + 2hx - 3 = 0.$ By \href{https://artofproblemsolving.com/wiki/index.php?title=Vieta\%27s\_Formulas}{Vieta's Formulas}, the sum of the roots is $-2h$ and the product of the roots is $-3.$

 Let the two roots be $r$ and $s.$ Note that
 \[r^2 + s^2 = (r+s)^2 - 2rs = (-2h)^2 -2(-3)\]

 Therefore, $4h^2 + 6 = 10$ and $h = \pm 1.$ This doesn't match any of the answer choices, so the answer is $\boxed{\textbf{(E) }\text{None of these}}.$

 -edited by coolmath34

 



\textbf{Solution 2}

The given equation can be rewritten as $x^2+2hx-3=0$. By \href{https://artofproblemsolving.com/wiki/index.php?title=Vieta\%27s\_Formulas}{Vieta's Formulas}, we know that the sum of the roots is $-2h$. Thus, by \href{https://artofproblemsolving.com/wiki/index.php?title=Newton\%27s\_Sums}{Newton's sums}, we have the following equation:\begin{align*}10+(2h)(-2h)+2(-3) &= 0 \\10-4h^2-6 &= 0 \\-4h^2 &= -4 \\h^2 &= 1 \\|h| &= 1 \\\end{align*}
Thus, our answer is $\boxed{\textbf{(E) }\text{None of these}}$.

 


	\item (\textit{AHSME 1971}) If $\log_2(\log_3(\log_4 x))=\log_3(\log_4(\log_2 y))=\log_4(\log_2(\log_3 z))=0$, then the sum $x+y+z$ is equal to


 $\textbf{(A) }50\qquad \textbf{(B) }58\qquad \textbf{(C) }89\qquad \textbf{(D) }111\qquad  \textbf{(E) }1296$


 
 \textbf{Answer: }C



\textbf{Solution}

If $\log_{x}{n}=0,$ then $n = 1.$ So, we can rewrite this equation:
 \[(\log_3(\log_4 x))=(\log_4(\log_2 y))=(\log_2(\log_3 z))=1\]

 Solve individually for each variable.
 \[x = 4^3 = 64\]
 \[y = 2^4 = 16\]
 \[z = 3^2 = 9\]Therefore, $x + y + z = 89.$

 The answer is $\boxed{\textbf{(C) }89}.$

 -edited by coolmath34

 


	\item (\textit{AHSME 1971}) If $w$ is one of the imaginary roots of the equation $x^3=1$, then the product $(1-w+w^2)(1+w-w^2)$ is equal to


 $\textbf{(A) }4\qquad \textbf{(B) }w\qquad \textbf{(C) }2\qquad \textbf{(D) }w^2\qquad  \textbf{(E) }1$


 
 \textbf{Answer: }A



\textbf{Solution 1}

Expanding the given expression yields $1+w-w^2-w-w^2+w^3+w^2+w^3-w^4=1-w^2+2w^3-w^4$. Recalling that $w^3=1$, we see that this expression equals $1-w^2+2-w=4-(1+w+w^2)$. By the properties of \href{https://artofproblemsolving.com/wiki/index.php?title=Roots\_of\_unity}{roots of unity} $\neq 1$, $w^2+w+1=0$, so the given expression equals $\boxed{\textbf{(A) }4}$.

 



\textbf{Solution 2 (not recommended)}

Suppose $w=e^{i\frac{2\pi}3}=\cos(\frac{2\pi}3)+i\sin(\frac{2\pi}3) = \frac{-1+i\sqrt3}2$. Substituting this into the given expression, we can calculate the result:\begin{align*}(1-w+w^2)(1+w-w^2) &= (1-\frac{-1+i\sqrt3}2+(\frac{-1+i\sqrt3}2)^2)(1+\frac{-1+i\sqrt3}2-(\frac{-1+i\sqrt3}2)^2) \\&= (1-\frac{-1+i\sqrt3}{2}+\frac{1-3-2i\sqrt3}{4})(1+\frac{-1+i\sqrt3}2-\frac{1-3-2i\sqrt3}{4}) \\&= (1-\frac{-1+i\sqrt3}{2}+\frac{-1-i\sqrt3}{2})(1+\frac{-1+i\sqrt3}2-\frac{-1-i\sqrt3}{2}) \\&= (1-2(\frac{i\sqrt3}2))(1+2(\frac{i\sqrt3}2)) \\&= 1^2-(i\sqrt3)^2 \\&= 1+3 \\&= 4\end{align*}
Thus, our answer is $\boxed{\textbf{(A) }4}$.

 


	\item (\textit{AHSME 1971}) Teams $A$ and $B$ are playing a series of games. If the odds for either to win any game are even and 
Team $A$ must win two or Team $B$ three games to win the series, then the odds favoring Team $A$ to win the series are


 $\textbf{(A) }11\text{ to }5\qquad \textbf{(B) }5\text{ to }2\qquad \textbf{(C) }8\text{ to }3\qquad \textbf{(D) }3\text{ to }2\qquad \textbf{(E) }13\text{ to }6$


 
 \textbf{Answer: }A



\textbf{Solution}

We have two cases: one where $A$ wins the first game and the other where $A$ loses.

 $\underline{\text{Case 1:}}$ In the $\frac12$ chance that $A$ wins the first game, $A$ simply needs to win at least $1$ of the next $3$ games. We see that the probability of $A$ losing the next $3$ games is $(\frac12)^3=\frac18$, so, by \href{https://artofproblemsolving.com/wiki/index.php?title=Complementary\_counting}{complementary counting}, the probability that $A$ wins at least $1$ of the next $3$ games is $1-\frac18=\frac78$.

 $\underline{\text{Case 2:}}$ In the $\frac12$ chance that $A$ loses the first game, both teams need to win $2$ games, so $A$'s advantage completely disappears. Thus, the proability that $A$ wins the series from here is $\frac12$.

 Combining the information from the two cases, we see that $A$'s probability of winning the series is $\frac12 \cdot \frac78 + \frac12 \cdot \frac12 = \frac7{16} + \frac4{16} = \frac{11}{16}$. Thus, our answer is $\boxed{\textbf{(A) } 11 \text{ to } 5}$.

 


	\item (\textit{AHSME 1971}) 
                    \begin{figure}[H]
                        centering
                        \includegraphics[width=0.5\textwidth]{Problem/AHSME_I_1971_24_Problem_Figure_1.png}
                    \end{figure}
                    


 Pascal's triangle is an array of positive integers(See figure), in which the first row is $1$, the second row is two $1$'s, 
each row begins and ends with $1$, and the $k^\text{th}$ number in any row when it is not $1$, is the sum of the $k^\text{th}$
and $(k-1)^\text{th}$ numbers in the immediately preceding row. The quotient of the number of numbers in the first $n$ rows 
which are not $1$'s and the number of $1$'s is


 $\textbf{(A) }\dfrac{n^2-n}{2n-1}\qquad \textbf{(B) }\dfrac{n^2-n}{4n-2}\qquad \textbf{(C) }\dfrac{n^2-2n}{2n-1}\qquad \textbf{(D) }\frac{n^2-3n+2}{4n-2}\qquad \textbf{(E) }\text{None of these}$


 
 \textbf{Answer: }D



\textbf{Solution}

The first row has $1$ one, and all the other rows have $2$ ones. Thus, the total number of ones in the first $n$ rows is $1+2(n-1)=2n-1$. The total number of elements in the first $n$ rows is $1+2+3+\ldots+n$, which is the $n$th \href{https://artofproblemsolving.com/wiki/index.php?title=Triangular\_number}{triangular number}, $\frac{n(n+1)}2$. Thus, the number of elements excluding ones is $\frac{n(n+1)}2-(2n-1)$. Thus, our desired quotient is $\frac{\frac{n(n+1)}2-(2n-1)}{2n-1} = \frac{n^2+n-2(2n-1)}{2(2n-1)} = \boxed{\textbf{(D) }\frac{n^2-3n+2}{4n-2}}$.

 


	\item (\textit{AHSME 1971}) A teen age boy wrote his own age after his father's. From this new four place number, he subtracted the absolute value of the 
difference of their ages to get $4,289$. The sum of their ages was


 $\textbf{(A) }48\qquad \textbf{(B) }52\qquad \textbf{(C) }56\qquad \textbf{(D) }59\qquad  \textbf{(E) }64$


 
 \textbf{Answer: }D



\textbf{Solution}

Because the father's age is a two digit number, we know that the father's age must be $43$ so that the second two digits of the original number can be a number in the teens. Let the son's age be $x$. We know that the original number is $4300+x$, and the positive difference between their ages is $43-x$. Thus, we have the equation $4300+x-(43-x)=4289$, which yields $x=16$. Thus, the sum of the two ages is $43+16=\boxed{\textbf{(D) }59}$.

 


	\item (\textit{AHSME 1971}) 
                    \begin{figure}[H]
                        centering
                        \includegraphics[width=0.5\textwidth]{Problem/AHSME_I_1971_26_Problem_Figure_1.png}
                    \end{figure}
                    


 In $\triangle ABC$, point $F$ divides side $AC$ in the ratio $1:2$. Let $E$ be the point of intersection of side $BC$ and $AG$ where $G$ is the 
midpoint of $BF$. The point $E$ divides side $BC$ in the ratio


 $\textbf{(A) }1:4\qquad \textbf{(B) }1:3\qquad \textbf{(C) }2:5\qquad \textbf{(D) }4:11\qquad  \textbf{(E) }3:8$


 
 \textbf{Answer: }B



\textbf{Solution 1}

We will use \href{https://artofproblemsolving.com/wiki/index.php?title=Mass\_points}{mass points} to solve this problem. $AC$ is in the ratio $1:2,$ so we will assign a mass of $2$ to point $A,$ a mass of $1$ to point $C,$ and a mass of $3$ to point $F.$

 We also know that $G$ is the midpoint of $BF,$ so $BG:GF=1:1.$ $F$ has a mass of $3,$ so $B$ also has a mass of $3.$

 In line $BC,$ $B$ has a mass of $3$ and $C$ has a mass of $1.$ Therefore, $BE:EC = 1:3.$

 The answer is $\boxed{\textbf{(B) }1:3}.$

 -edited by coolmath34

 



\textbf{Solution 2}

By \href{https://artofproblemsolving.com/wiki/index.php?title=Menelaus\%27\_Theorem}{Menelaus' Theorem} on $\triangle AGF$ and line $\overleftrightarrow{BC}$, we know that $\frac{AE}{GE} \cdot \frac{GB}{FB} \cdot \frac{FC}{AC}=1$. Because $G$ is the midpoint of $\overline{FB}$, we know that $\frac{GB}{FB}=\frac12$. Furthermore, from the problem, we know $\frac{FC}{AC}=\frac23$, so, by susbtitution,  our first equation becomes $\frac{AE}{GE} \cdot \frac12 \cdot \frac23 = 1$, so $\frac{AE}{GE}=3$, and therefore $\frac{AE}{AG}=\frac32$.

 Using Menelaus again on $\triangle BEG$ and line $\overleftrightarrow{AC}$ reveals that $\frac{BC}{EC} \cdot \frac{EA}{GA} \cdot \frac{GF}{BF}=1$. We know that $\frac{EA}{GA}=\frac{AE}{AG}=\frac32$ and $\frac{GF}{BF}=\frac12$. Thus, our equation becomes $\frac{BC}{EC} \cdot \frac32 \cdot \frac12=1$, and so  $\frac{BC}{EC}=\frac43$. Because $\frac{BC}{EC}=\frac{BE+EC}{EC}=1+\frac{BE}{EC}$, we see that $\frac{BE}{EC}=\frac13$, so our answer is $\boxed{\textbf{(B) }1:3}$.

 


	\item (\textit{AHSME 1971}) A box contains chips, each of which is red, white, or blue. The number of blue chips is at least half the number of white chips, 
and at most one third the number of red chips. The number which are white or blue is at least $55$. The minimum number of red chips is


 $\textbf{(A) }24\qquad \textbf{(B) }33\qquad \textbf{(C) }45\qquad \textbf{(D) }54\qquad  \textbf{(E) }57$


 
 \textbf{Answer: }E



\textbf{Solution}

Let the number of white be $2x$. The number of blue is then $x-y$ for some constant $y$. So we want $2x+x-y=55\rightarrow 3x-y=55$. We take mod 3 to find y. $55=1\pmod{3}$, so blue is 1 more than half white. The number of whites is then 36, and the number of blues is 19. So $19*3=\boxed{\textbf{(E) }57}$

 


	\item (\textit{AHSME 1971}) Nine lines parallel to the base of a triangle divide the other sides each into $10$ equal segments and the area into $10$ distinct parts. 
If the area of the largest of these parts is $38$, then the area of the original triangle is 


 $\textbf{(A) }180\qquad \textbf{(B) }190\qquad \textbf{(C) }200\qquad \textbf{(D) }210\qquad  \textbf{(E) }240$


 
 \textbf{Answer: }C



\textbf{Solution}


                    \begin{figure}[H]
                        centering
                        \includegraphics[width=0.5\textwidth]{Solution/AHSME_I_1971_28_Solution_Figure_1.png}
                    \end{figure}
                    

 Let the triangle be $\triangle ABC$ with base $\overline{BC}$ and longest parallel segment $\overline{DE}$ with $D$ on $\overline{AB}$ and $E$ on $\overline{AC}$, as in the diagram.

 By the properties of \href{https://artofproblemsolving.com/wiki/index.php?title=Transversals}{transversals}, we have $\measuredangle ABC = \measuredangle ADE$. Thus, by \href{https://artofproblemsolving.com/wiki/index.php?title=AA\_Similarity}{AA Similarity}, we have $\triangle ADE \sim \triangle ABC$ (because they share $\angle BAC$). From the problem, we know that $\frac{AE}{AC}=\frac9{10}$, so, by similarity, $\frac{DE}{BC}=\frac9{10}$, and so $DE=\frac9{10}BC$.

 Now, let $[\triangle ABC] = A$. Because $\frac{DE}{BE}=\frac9{10}$, we know that $\frac{[\triangle ADE]}{[\triangle ABC]}=(\frac9{10})^2=\frac{81}{100}$. From the problem, $[BCED]=38$, so $A=[\triangle ADE]+38=\frac9{10}A+38$. Solving for $A$ yields  $A=\boxed{\textbf{(C) }200}$.

 


	\item (\textit{AHSME 1971}) Given the progression $10^{\dfrac{1}{11}}, 10^{\dfrac{2}{11}}, 10^{\dfrac{3}{11}}, 10^{\dfrac{4}{11}},\dots , 10^{\dfrac{n}{11}}$. 
The least positive integer $n$ such that the product of the first $n$ terms of the progression exceeds $100,000$ is


 $\textbf{(A) }7\qquad \textbf{(B) }8\qquad \textbf{(C) }9\qquad \textbf{(D) }10\qquad  \textbf{(E) }11$


 
 \textbf{Answer: }E



\textbf{Solution}

By the rules of \href{https://artofproblemsolving.com/wiki/index.php?title=Exponentiation}{exponents}, the product of the first $n$ terms of the sequence equals $10^{\frac{1}{11}+\frac{2}{11}+\dots+\frac{n}{11}}$. From here, we can set up the equation $10^{\frac{1}{11}+\frac{2}{11}+\dots+\frac{n}{11}}=100,000=10^5$, which simplifies to $\frac{1}{11}+\frac{2}{11}+\dots+\frac{n}{11}=5$, or $1+2+3\dots n=55$. By the formula for \href{https://artofproblemsolving.com/wiki/index.php?title=Triangular\_numbers}{triangular numbers}, $\frac{n(n+1)}{2}=55$. Thus, $n^2+n-110=0$, or $(n+11)(n-10)=0$. Because $n>0$, only $n=10$ makes the product equal $100,000$. However, we have to strictly exceed $100,000$, so our answer is $\boxed{\textbf{(E) }11}.$

 


	\item (\textit{AHSME 1971}) Given the linear fractional transformation of $x$ into $f_1(x)=\dfrac{2x-1}{x+1}$. Define $f_{n+1}(x)=f_1(f_n(x))$ for $n=1,2,3,\cdots$. 
Assuming that $f_{35}(x)=f_5(x)$, it follows that $f_{28}(x)$ is equal to


 $\textbf{(A) }x\qquad \textbf{(B) }\frac{1}{x}\qquad \textbf{(C) }\frac{x-1}{x}\qquad \textbf{(D) }\frac{1}{1-x}\qquad  \textbf{(E) }\text{None of these}$


 
 \textbf{Answer: }D



\textbf{Solution 1}

Extend the definition of $f_n(x)$ to $n\leq0$: Let $f_{n-1}(x)$ be the function such that $f_1(f_{n-1}(x))=f_n(x)$. From the problem, $f_5(x)=f_{35}(x)$, so the functions $f_1(x),f_2(x),\ldots$ must repeat in a cycle whose length is a cycle which is a divisor of $35-5=30$. Thus, if $f_a(x)=f_b(x)$ for integers $a$ and $b$, we know that $a\equiv b$ \href{https://artofproblemsolving.com/wiki/index.php?title=Modulo}{modulo} $30$. Thus, because $28\equiv-2\pmod{30}$, we know that $f_{-2}(x)=f_{28}(x)$.

 It is clear that $f_0(x)=x$, because $f_1(f_0(x))=f_1(x)$.

 Let $f_{-1}(x)=y$. Then, we know that $f_1(y)=f_0(x)=x$, so we have the following equation we can solve for $y$:\begin{align*}\frac{2y-1}{y+1} &= x \\2y-1 &= xy+x \\y(2-x) &= x+1 \\y &= \frac{x+1}{2-x}\end{align*}


 Let $f_{-2}(x)=z$. Then, we know that $f_1(z)=f_{-1}(x)=\frac{x+1}{2-x}$, so we have the following equation we can solve for $z$:\begin{align*}\frac{2z-1}{z+1} &= \frac{x+1}{2-x} \\4z-2-2xz+x &= xz+x+z+1 \\3z-3xz &= 3 \\z(1-x) &= 1 \\z &= \frac1{1-x}\end{align*}


 We derived earlier the fact that $z=f_{-2}(x)=f_{28}(x)$, so $f_{28}(x)=\boxed{\textbf{(D) }\frac{1}{1-x}}$.

 



\textbf{Solution 2}

Keep solving for the next function in the sequence of $f_n(x)$ while being sure not to make silly algebra mistakes. This process reveals that $f_6(x)=x$, so $f_7(x)=f_1(x)$, and the cycle of functions repeats \href{https://artofproblemsolving.com/wiki/index.php?title=Modulo}{modulo} $6$. Because $28\equiv4\pmod6$, we know that $f_{28}(x)=f_4(x)=\boxed{\textbf{(D) }\frac{1}{1-x}}$, which we calculated on the way to deducing that $f_6(x)=x$.

 


	\item (\textit{AHSME 1971}) 
                    \begin{figure}[H]
                        centering
                        \includegraphics[width=0.5\textwidth]{Problem/AHSME_I_1971_31_Problem_Figure_1.png}
                    \end{figure}
                    


 Quadrilateral $ABCD$ is inscribed in a circle with side $AD$, a diameter of length $4$. If sides $AB$ and $BC$ each have length $1$, 
then side $CD$ has length


 $\textbf{(A) }\frac{7}{2}\qquad \textbf{(B) }\frac{5\sqrt{2}}{2}\qquad \textbf{(C) }\sqrt{11}\qquad \textbf{(D) }\sqrt{13}\qquad \textbf{(E) }2\sqrt{3}$


 
 \textbf{Answer: }A



\textbf{Solution 1}

Note that the length 4 forms a semicircle. We can then use the Law of Cosines. Take the center and form a line segment with the other two points.

 Let's find the cosine of angle $AOB$. The cosine of that angle is 7/8. We use the double cosine angle to find angle $AOC$: 

 $cos(2x) = 2 cos^2 (x)- 1$

 Therefore, $cos(2x) = 17/32$

 Applying the Law of Cosines to the triangle $DOC$, we get 

 $CD^2 = 2^2 + 2^2 - 2(4)(-17/32)$ (because $cos(x) = -cos(180-x)$)

 Therefore, $CD^2 = 8 + 17/4$

 and we find that $CD=\boxed{\textbf{(A) }7/2}$!

 



\textbf{Solution 2}

By the properties of a circle, we see that angle $\text{DCA}$ and angle $\text{DBA}$ are both right angles. Thus, if we formulate the diagonal lengths of the cyclic quadrilateral using the Pythagorean theorem, we can finish the problem with Ptolemy's Theorem. 

 With that train of thought, we see that 

 
 \[BD^2=AD^2-AB^2\]

 
 \[BD^2=4^2-1^2\]

 
 \[BD^2=15 \Rightarrow BD = \sqrt{15}\]

 To let us formulate that rest of our lengths, let $CD=x$. Then, similar to the above, 

 
 \[AC^2=AD^2-CD^2\]

 
 \[AC^2=4^2-x^2=16-x^2\]

 
 \[AC=\sqrt{16-x^2}\]

 Now that we have all of the side lengths and diagonal lengths in one variable, we use Ptolemy's Theorem to finish from here: 

 
 \[4\cdot1 + x\cdot1 = \sqrt{15}\sqrt{16-x^2}\]

 
 \[x+4=\sqrt{240-15x^2}\]

 
 \[x^2+8x+16=240-15x^2\]

 
 \[16x^2+8x-224=0\]

 
 \[2x^2+x-28=0\]

 
 \[x=\frac{-1 \pm \sqrt{1-(4)(2)(-28)}}{2(2)}=\frac{-1\pm \sqrt{225}}{4}=\frac{-1 \pm 15}{4}\]

 Since the negative solution is extraneous, we see that our solution is $x=\frac{7}{2}\Rightarrow\boxed{\textbf{(A)}}$

 


	\item (\textit{AHSME 1971}) If $s=(1+2^{-\frac{1}{32}})(1+2^{-\frac{1}{16}})(1+2^{-\frac{1}{8}})(1+2^{-\frac{1}{4}})(1+2^{-\frac{1}{2}})$, then $s$ is equal to


 $\textbf{(A) }\textstyle{\frac{1}{2}}(1-2^{-\frac{1}{32}})^{-1}\qquad \textbf{(B) }(1-2^{-\frac{1}{32}})^{-1}\qquad \textbf{(C) }1-2^{-\frac{1}{32}}\qquad \\ \textbf{(D) }\textstyle{\frac{1}{2}}(1-2^{-\frac{1}{32}})\qquad  \textbf{(E) }\frac{1}{2}$


 
 \textbf{Answer: }A



\textbf{Solution 1}

Multiply both sides of the given equation by $(1-2^{-\frac1{32}})$. Using \href{https://artofproblemsolving.com/wiki/index.php?title=Difference\_of\_squares}{difference of squares} reveals that $(1-2^{-\frac1{32}})(1+2^{-\frac1{32}})=(1-2^{-\frac1{16}})$. Thus, we can use difference of squares again on the next term, $(1+2^{-\frac1{16}})$, and the terms after that until we get $(1+2^{-\frac1{32}})s=(1-2^{-\frac12})(1+2^{-\frac12})=1-2^{-1}=\frac12$. Dividing both sides by $(1+2^{-\frac1{32}})$ reveals that $s=\boxed{\textbf{(A) }\frac12(1+2^{-\frac1{32}})^{-1}}$.

 



\textbf{Solution 2 (50-50 guess)}

For any two reals $a$ and $x$ with $a>0$, $a^x>0$. Thus, we know that all of the terms $2^{-\frac1x}>0$, so $s$ is the product of $5$ real terms greater than $1$. Thus, $s>1$, so we can eliminate options (C), (D), and (E). Now, we have a $50$\% chance of guessing the right answer, $\boxed{\textbf{(A) }\frac12(1+2^{-\frac1{32}})^{-1}}$.

 


	\item (\textit{AHSME 1971}) If $P$ is the product of $n$ quantities in Geometric Progression, $S$ their sum, and $S'$ the sum of their reciprocals, 
then $P$ in terms of $S, S'$, and $n$ is


 $\textbf{(A) }(SS')^{\frac{1}{2}n}\qquad \textbf{(B) }(S/S')^{\frac{1}{2}n}\qquad \textbf{(C) }(SS')^{n-2}\qquad \textbf{(D) }(S/S')^n\qquad \textbf{(E) }(S/S')^{\frac{1}{2}(n-1)}$


 
 \textbf{Answer: }B



\textbf{Solution 1}

Let the \href{https://artofproblemsolving.com/wiki/index.php?title=Geometric\_sequence}{geometric sequence} have first term $a$ and common ratio $R$. Then, the first $n$ terms of the sequence are $a,aR,aR^2,\ldots,aR^{n-1}$. The product of these terms $P$ is $a^nR^{1+2+\ldots+n-1}=a^nR^{\frac{(n-1)n}2}$ by the formula for \href{https://artofproblemsolving.com/wiki/index.php?title=Triangular\_numbers}{triangular numbers}. Using the \href{https://artofproblemsolving.com/wiki/index.php?title=Geometric\_sequence\#Sum}{sum formula} reveals that $S=a\cdot\frac{1-R^n}{1-R}$.

 We know that $S^{\prime}=\frac1a+\frac1{aR}+\ldots+\frac1{aR^{n-1}}$ Combining fractions reveals that $S^{\prime}=\frac{aR^{n-1}+aR^{n-2}+\ldots+aR+a}{a^2R^{n-1}}=\frac S{a^2R^{n-1}}$. Note that this denominator looks suspiciously similar to our formula for $P$. In fact, $(a^2+R^{n-1})^{\frac n2}=a^n+R^{\frac{(n-1)n}2}=P$. Because $(S/S^{\prime})^{\frac n2}=(a^2+R^{n-1})^{\frac n2}=P$, our answer is $\boxed{\textbf{(B) }(S/S^{\prime})^{\frac 12n}}$.

 



\textbf{Solution 2 (Answer Choices)}

We can just look at a very specific case: $1, 2, 4, 8.$ Here, $n=4, P=64, S=15,$ and $S'=\frac{30}{16}=\frac{15}{8}.$

 Then, plug in values of $S, S',$ and $n$ into each of the answer choices and see if it matches the product. 

 Answer choice $\boxed{\textbf{(B) }(S/S')^{\frac{1}{2}n}}$ works: $(\frac{15}{\frac{15}{8}})^2 = 64.$

 -edited by coolmath34

 



\textbf{Solution 3 (Answer Choices)}

We can use \href{https://artofproblemsolving.com/wiki/index.php?title=Dimensional\_analysis&action=edit&redlink=1}{dimensional analysis} to cut down our answer choices. Suppose that each of the terms in the geometric progression is in units of $\text{meters}$. Then, $S$ should have units of $\text{meters}$, $S^{\prime}$ units of $\frac1{\text{meters}}$ and $P$ units of $\text{meters}^n$. Therefore, $SS^{\prime}$ is unitless, so we can eliminate options (A) and (C). $S/S^{\prime}$ has units $\text{meters}^2$, so, to equal $P$ (which has units $\text{meters}^n$), the exponent needs to be $\frac n2$. The only remaining answer choice which satsifies this constraint is $\boxed{\textbf{(B) }(S/S^{\prime})^{\frac12n}}$.

 


	\item (\textit{AHSME 1971}) An ordinary clock in a factory is running slow so that the minute hand passes the hour hand at the usual dial position($12$ o'clock, etc.) 
but only every $69$ minutes. At time and one-half for overtime, the extra pay to which a $\textdollar 4.00$ per hour worker should be entitled 
after working a normal $8$ hour day by that slow running clock, is


 $\textbf{(A) }\textdollar 2.30\qquad \textbf{(B) }\textdollar 2.60\qquad \textbf{(C) }\textdollar 2.80\qquad \textbf{(D) }\textdollar 3.00\qquad \textbf{(E) }\textdollar 3.30$


 
 \textbf{Answer: }B



\textbf{Solution}

If the clock were functioning properly, the minute hand would have an angular velocity of $1\frac{\text{rev}}{\text{hr}}$, and the hour hand would have an angular velocity of $\frac1{12}\frac{\text{rev}}{\text{hr}}$. Thus, the minute hand would move at a velocity of $1-\frac1{12}=\frac{11}{12}\frac{\text{rev}}{\text{hr}}$ relative to the hour hand. Thus, the minute hand would pass the hour hand every $\frac{12}{11}$ hours, or $\frac{720}{11}$ minutes (after the minute hand has made a full revolution relative to the hour hand). Because the slow clock takes $69$ minutes to do this, one hour to this clock is actually $\frac{69}{(\frac{720}{11})}=\frac{253}{240}$ hours. Thus, $8$ hours on the slow clock is actually $8\cdot\frac{253}{240}=\frac{253}{30}$ hours. Thus, the workers worked for an extra $\frac{253}{30}-8=\frac{13}{30}$ hours. Because, from the problem, the workers are paid $1.5$ times their usual wage when working overtime, they earn $\$4\cdot1.5=\$6$ per hour in overtime. Therefore, they should receive $\frac{13}{30}\cdot6=\frac{13}{5}=2.6$ extra dollars, so our answer is $\boxed{\textbf{(B) }\$2.60}$.

 


	\item (\textit{AHSME 1971}) Each circle in an infinite sequence with decreasing radii is tangent externally to the one following it and to both sides of a given right angle. 
The ratio of the area of the first circle to the sum of areas of all other circles in the sequence, is


 $\textbf{(A) }(4+3\sqrt{2}):4\qquad \textbf{(B) }9\sqrt{2}:2\qquad \textbf{(C) }(16+12\sqrt{2}):1\qquad \\ \textbf{(D) }(2+2\sqrt{2}):1\qquad  \textbf{(E) }(3+2\sqrt{2}):1$


 
 \textbf{Answer: }C



\textbf{Solution 1}


                    \begin{figure}[H]
                        centering
                        \includegraphics[width=0.5\textwidth]{Solution/AHSME_I_1971_35_Solution_Figure_1.png}
                    \end{figure}
                    

 \href{https://artofproblemsolving.com/wiki/index.php?title=WLOG}{WLOG}, let the radius of the largest circle be $1$ (we can do this because scaling the diagram preserves the ratio of areas). Let the largest circle be $\omega_1$ and the second largest circle be $\omega_2$. As in the diagram above, let $\omega_1$ have center $A$ with points of tangency at $B$ and $D$. Let the tangents from these two points meet at $C$. From the problem, we know that $\measuredangle BCD=90^{\circ}$. Let the common tangent of $\omega_1$ and $\omega_2$ intersect $\overline{BC}$ and $\overline{DC}$ at $X$ and $Y$, respectively. Let $\omega_2$ intersect $\overline{XY}$ at $R$.

 By the \href{https://artofproblemsolving.com/wiki/index.php?title=Two\_Tangent\_Theorem}{Two Tangent Theorem}, $CB=CD$. Thus, $\triangle BCD$ is an isosceles right triangle, so $\measuredangle CBD=\measuredangle CDB=45^{\circ}$. $\measuredangle CBD$ is a \href{https://artofproblemsolving.com/wiki/index.php?title=Tangent-chord\_angle&action=edit&redlink=1}{tangent-chord angle} in $\omega_1$, so $\measuredangle BAD=2\cdot\measuredangle CDB=90^{\circ}$. Thus, because $ABCD$ has $4$ right angles and $2$ adjacent sides congruent, it is a square. Thus, $BC=CD=1$.

 $\omega_2$ is the \href{https://artofproblemsolving.com/wiki/index.php?title=Incircle}{incircle} of $\triangle CXY$, because it is tangent to all its sides. Let $CX=x$ and $s$ be the \href{https://artofproblemsolving.com/wiki/index.php?title=Semiperimeter}{semiperimeter} of $\triangle CXY$. Because $\overline{XY}\parallel\overline{BD}$, we see that $\measuredangle CXY=\measuredangle CBD=45^{\circ}$ and $\measuredangle CYX=\measuredangle CDB=45^{\circ}$. Thus, $\triangle CXY$ is an isosceles right triangle, so $XY=x\sqrt2$. By the properties of an incircle's touchpoints, $XR=s-CY=s-x$ and $YR=s-CX=s-x$. Thus, $R$ is the midpoint of $\overline{XY}$, so $RX=XY=\frac{\sqrt2}2x$.

 Because $\overline{XY}$ is the common tangent of $\omega_1$ and $\omega_2$, $R$ is on $\omega_1$. Thus, by the Two Tangent Theorem, $XB=XR=\frac{\sqrt2}2x$.

 Now, we have the equation $1=BC=BX+XC=\frac{\sqrt2}2x+x$, which yields $x=\frac1{1+\frac{\sqrt2}2}=\frac{2(2-\sqrt2)}{(2+\sqrt2)(2-\sqrt2)}=2-\sqrt2$. We know that $s=\frac{x+x+x\sqrt2}{2}=x+\frac{\sqrt2}2x$. Let the \href{https://artofproblemsolving.com/wiki/index.php?title=Inradius}{inradius} of $\triangle CXY$ be $r$. Now the formula $A=rs$ in $\triangle CXY$ gives us the equation $\frac{x^2}2=r(x+\frac{\sqrt2}2x)$, which yields $r=\frac x{(2+\sqrt2)}=\frac{(2-\sqrt2)^2}{(2+\sqrt2)(2-\sqrt2)}=\frac{4-4\sqrt2+2}{4-2}=3-2\sqrt2$. 

 Because the ratio of the radii of $\omega_1$ and $\omega_2$ is $\frac{3-2\sqrt{2}}1=3-2\sqrt2$, the ratio of their areas is $(3-2\sqrt2)^2=9-12\sqrt2+8=17-12\sqrt2$. This ratio is the same as that of the area of $\omega_2$ to the third largest circle, the third to the fourth, etc., because removing the largest circle gives a scaled-down version of the same problem. Thus,  the total area of the circles is given by the \href{https://artofproblemsolving.com/wiki/index.php?title=Geometric\_sequence}{geometric sequence} starting at $\pi$ (the area of $\omega_1$) with common ratio $17-12\sqrt2\approx17-12\cdot1.4=0.2<1$, which, by the infinite \href{https://artofproblemsolving.com/wiki/index.php?title=Geometric\_series}{geometric series} formula, has an infinite sum of $\pi\cdot\frac1{1-(17-12\sqrt2)}=\pi\cdot\frac{(12\sqrt2+16)}{(12\sqrt2-16)(12\sqrt2+16)}=\pi\cdot\frac{3\sqrt2+4}8$. However, we are looking for the ratio of the area of $\omega_1$ to that of all of the circles \textit{excluding} $\omega_1$, so we have to subtract $\pi$. This reveals that the area of the other circles is $\pi\cdot\frac{3\sqrt2-4}8$. Now, our desired answer is the ratio $\frac{\pi}{\pi\cdot\frac{3\sqrt2-4}8}=\frac{8(3\sqrt2+4)}{(3\sqrt2-4)(3\sqrt2+4)}=\frac{8(3\sqrt2+4)}{18-16}=4(3\sqrt2+4)=12\sqrt2+16$. Thus, we choose $\boxed{\textbf{(C) }(16+12\sqrt2):1}$.

 



\textbf{Solution 2 (Informed guess)}

Draw a good diagram, and especially make sure that you are drawing a right angle. Then, you will realize that the big circle completely dwarfs all of the other circles. Answer choice (C), $16+12\sqrt2$, is about $16+12\cdot1.4=32.8$, whereas the greatest of the other answer choices is about $6.3$, which is unreasonably small based on our diagram. Thus, we choose answer $\boxed{\textbf{(C) }(16+12\sqrt2):1}$.

 


\end{enumerate}